\documentclass[10pt,a4paper]{amsart}
\usepackage[margin=19mm]{geometry}
\usepackage{amsmath,amssymb,mathtools}
\usepackage[scr=rsfs,cal=euler]{mathalfa}
\usepackage{xfrac}
\usepackage[unicode,hidelinks,bookmarks=false]{hyperref}
\newcommand{\ie}{i.e.\ }
\newcommand{\cf}{cf.\ }
\newcommand{\resp}{respectively\ }
\newcommand{\Subsection}{\subsection}
\providecommand{\nobreakdash}{\nobreak\hbox{-}}
\setlength{\emergencystretch}{2em}
\multlinegap0pt
\usepackage{mathtools}
\usepackage{amssymb}
\usepackage{bm}
\newtheorem{lemma}{Lemma}
\newtheorem{proposition}{Proposition}
\title{Homotopy commutativity in quasitoric manifolds}
\author{Sho Hasui, Daisuke Kishimoto, Yichen Tong, Mitsunobu Tsutaya}
\date{Source: arXiv:2404.01510v2 (CC BY 4.0)}
\begin{document}
\maketitle

\noindent\emph{Source and licence.} Homotopy commutativity in quasitoric manifolds. Version arXiv:2404.01510v2 (CC BY 4.0), distributed by arXiv under the Creative Commons Attribution 4.0 International licence, http://creativecommons.org/licenses/by/4.0/, as recorded in the arXiv licence metadata for this exact version and shown on the abstract page https://arxiv.org/abs/2404.01510v2. Full licence notice: \texttt{licenses/arxiv-2404.01510-CC-BY-4.0.txt}.\par\smallskip

\noindent\textit{Editorial scope.} Counted unit: the article's proposition H(k,n) (source0.tex lines 841--844). The article's complete original proof of it is delivered with it (source0.tex lines 846--873, one \texttt{proof} environment of 28 lines). No mathematics is written, repaired or re-derived: the statement and the proof are the article's own lines, in the article's own notation, and no retained line is substituted or re-flowed.  Retained uncounted as the source-local context the proof needs: lines 818--825.  Nothing was omitted from any retained span, so the package carries no omission record.  No retained line contains a citation, so the wrapper carries no bibliography and no bibliography entry.  No retained line contains a figure environment, an image-inclusion command or drawing code, so no image is delivered and the package declares no asset dependency.  Editorial formatting only, in addition: an AMS \texttt{amsart} preamble that restates the article's own declarations and macros used by the retained lines, namely the environments \texttt{lemma}, \texttt{proposition} on separate counters, together with the standard packages those lines need.  The retained environments print in this document as Lemma 1, Proposition 1 in order of appearance, whereas the article prints the same items with the numbering of its own full document.  The complete native source of the article as distributed in the arXiv e-print for this exact version is bundled unchanged as \texttt{sources/157059/source0.tex}.
\begin{lemma}
  \label{I(k,n)}
  Every graded algebra isomorphism $f\colon H(k,n)\xrightarrow{\cong}H(l,n)$ satisfies
  \[
    f(I_j(k,n))=I_j(l,n)
  \]
  for $j=1,\ldots,n$.
\end{lemma}

\begin{proposition}
  \label{H(k,n)}
  $H(k,n)\cong H(l,n)$ if and only if $k=l$.
\end{proposition}

\begin{proof}
  The if part is trivial, and we consider the only if part. First, we consider the $n=2$ case. Suppose there is an isomorphism $f\colon H(k,2)\xrightarrow{\cong}H(l,2)$. By Lemma \ref{I(k,n)},
  \[
    f(t_1)=\epsilon_1(t_1+ct_2)\quad\text{and}\quad f(t_2)=\epsilon_2t_2
  \]
  for $\epsilon_1,\epsilon_2=\pm 1$ and an integer $c$, so we get
  \begin{align*}
    0&=f(t_1^4+kt_1^3t_2)\\
    &=(4c-l+k\epsilon_1\epsilon_2)t_1^3t_2+(6c^2+3k\epsilon_1\epsilon_2c)t_1^2t_2^2+(4c^3+3k\epsilon_1\epsilon_2c^2)t_1t_2^3.
  \end{align*}
  Then since $t_1^3t_2,\,t_1^2t_2^2,\,t_1t_2^3$ are linearly independent in $H(l,n)$, we obtain
  \[
    6c^2+3k\epsilon_1\epsilon_2c=0,\quad 4c^3+3k\epsilon_1\epsilon_2c^2=0,\quad 4c-l+k\epsilon_1\epsilon_2=0.
  \]
  By the first two equations, we get $c=0$, so by the third equation, we obtain $k=l$, as desired.



  Next, we consider the $n>2$ case. Let $H(m)$ denote the subalgebra of $H(m,n)$ generated by $t_{n-1}$ and $t_n$. Then there is a canonical isomorphism
  \[
    H(m)\cong H(m,2).
  \]
  Suppose there is an isomorphism $f\colon H(k,n)\xrightarrow{\cong}H(l,n)$. Then by Lemma \ref{I(k,n)}, the map $f$ restricts to an isomorphism
  \[
    H(k)\xrightarrow{\cong}H(l).
  \]
  Thus by the $n=2$ case, we get $k=l$, completing the proof.
\end{proof}

\end{document}
